\chapter{Fungsi Karakteristik dan Teorema Limit
Pusat}\label{ch:b3:clt}

Hukum bilangan besar mengatakan bahwa rata-ratanya konvergen; sedangkan teorema
limit pusat mengatakan \emph{bagaimana ia berfluktuasi}: bahwa galatnya,
yang diperbesar $\sqrt n$, bersifat \omterm{def:b3:clt:gaussianvector}{Gauss} secara asimtotik ---
apa pun distribusi asalnya. Keuniversalan ini adalah
fakta terdalam pada peluang elementer, dan bukti alaminya
bersifat Fourier: sebab \emph{\omterm{def:b3:clt:cf}{fungsi karakteristik}}
(yakni transformasi Fourier sebuah distribusi) mengubah jumlah yang bebas
menjadi hasil kali, lalu perkakas \cref{ch:b3:fouriertransform}
--- keinjektifan, titik tetap \omterm{def:b3:clt:gaussianvector}{Gauss} --- mengubah konvergensi
titik demi titik hasil kali itu menjadi konvergensi distribusinya
(yakni teorema Lévy, yang dibuktikan selengkapnya). Bab ini berakhir dengan
\omterm{def:b3:clt:gaussianvector}{vektor Gauss} dan penurunan yang jujur atas selang
kepercayaan yang dipakai di mana-mana dalam statistika; sedangkan soal akhir pekannya
menyajikan bukti kedua Lindeberg bagi teorema limit pusatnya, dengan laju
galat yang eksplisit.

\section{Fungsi karakteristik}

\begin{definition}\label{def:b3:clt:cf}
Yang disebut \emph{fungsi karakteristik}\index{fungsi
karakteristik} sebuah \omterm{def:b3:probability:space}{peubah acak} real $X$ adalah
\[
\varphi_X(\xi) = \E\bigl[\eu^{\iu\xi X}\bigr]
= \int_\R \eu^{\iu\xi x}\,\dd\P_X(x)
\qquad (\xi \in \R)
\]
(teorema transfernya menghitungnya dari distribusinya; sedangkan bagi \omterm{ex:b3:lebesgue:gamma}{kepadatan}
$f$, $\varphi_X(\xi) = \hat f(-\xi)$ dalam
konvensi \cref{ch:b3:fouriertransform}).
\end{definition}

\begin{proposition}\label{prop:b3:clt:cfbasics}
(a) Berlaku $\varphi_X(0) = 1$, $\abs{\varphi_X} \leq 1$, dan
$\varphi_X$ \omterm{def:b3:topology:continuity}{kontinu} seragam;
$\varphi_{aX + b}(\xi) = \eu^{\iu b\xi}\varphi_X(a\xi)$.
(b) Jika $X, Y$ \emph{bebas}:
$\varphi_{X+Y} = \varphi_X\,\varphi_Y$.
(c) Jika $\E\abs X^k < \infty$, maka $\varphi_X \in \mathcal
C^k$ dengan $\varphi_X^{(j)}(0) = \iu^j\,\E[X^j]$ bagi $j \leq
k$; khususnya, bagi $X \in L^2$ yang terpusat dan bervarians
$\sigma^2$:
\[
\varphi_X(\xi) = 1 - \frac{\sigma^2\xi^2}{2} +
o(\xi^2) \qquad (\xi \to 0).
\]
(d) \omterm{def:b3:clt:gaussianvector}{Gauss}: $X \sim \mathcal N(m, \sigma^2)$ mempunyai
$\varphi_X(\xi) = \eu^{\iu m\xi - \sigma^2\xi^2/2}$.
\end{proposition}

\begin{proof}
(a) Batasnya langsung; sedangkan untuk \omterm{def:b3:topology:continuity}{kekontinuannya}: $\abs{\varphi(\xi + h)
- \varphi(\xi)} \leq \E\abs{\eu^{\iu hX} - 1} \to 0$ saat $h
\to 0$ menurut konvergensi terdominasi, seragam terhadap $\xi$. Adapun
aturan afinnya adalah sebuah substitusi. (b) Berlaku $\eu^{\iu\xi(X+Y)} =
\eu^{\iu\xi X}\eu^{\iu\xi Y}$, sedangkan \omterm{def:b3:probability:space}{nilai harapan} hasil kali
variabel yang bebas memfaktor
(\cref{thm:b3:probability:independence}, yang diterapkan pada bagian real dan
imajinernya). (c) Penurunan di bawah tanda \omterm{def:b3:probability:space}{nilai harapannya},
yang terdominasi $\E\abs X^j$
(\cref{thm:b3:lebesgue:paramdiff}); lalu uraian Taylornya di
$0$ adalah Taylor--Young bagi fungsi $\mathcal C^2$
$\varphi$. (d) Untuk $\mathcal N(0,1)$: transformasi Gaussnya
(\cref{ex:b3:fouriertransform:gaussian} dengan $a = \frac12$)
memberikan $\int\eu^{\iu\xi x}\frac{\eu^{-x^2/2}}{\sqrt{2\pi}}\dd
x = \eu^{-\xi^2/2}$; sedangkan kasus umumnya lewat aturan afinnya.
\end{proof}

\begin{theorem}[Keinjektifan]\label{thm:b3:clt:injectivity}
Jika $\varphi_X = \varphi_Y$, maka $X$ dan $Y$ berdistribusi
sama. Lebih tepatnya, bagi $N \sim \mathcal N(0,1)$ yang bebas
dari $X$ dan $\varepsilon > 0$, variabel yang dimuluskan $X +
\varepsilon N$ mempunyai \omterm{ex:b3:lebesgue:gamma}{kepadatan}
\[
p_\varepsilon(x) = \frac1{2\pi}\int_\R
\varphi_X(-\xi)\,\eu^{-\varepsilon^2\xi^2/2}\,
\eu^{\iu\xi x}\,\dd\xi ,
\]
yang ditentukan $\varphi_X$ semata; lalu membiarkan $\varepsilon \to 0$
memulihkan distribusi $X$.\index{fungsi karakteristik,
keinjektifan}
\end{theorem}

\begin{proof}
Variabel $X + \varepsilon N$ berkepadatan $p_\varepsilon(x) =
\E\bigl[g_\varepsilon(x - X)\bigr]$, dengan $g_\varepsilon$ sebagai
\omterm{ex:b3:lebesgue:gamma}{kepadatan} $\mathcal N(0, \varepsilon^2)$: sebab bagi $B$ yang
Borel, \omterm{def:b3:probability:independence}{kebebasan} dan Tonelli memberikan $\P(X + \varepsilon N \in
B) = \int\!\!\int\mathbf 1_B(x + \varepsilon
n)g_1(n)\,\dd n\,\dd\P_X(x) = \int_B\E[g_\varepsilon(t -
X)]\dd t$ (substitusikan, lalu Tonelli lagi). Lalu dengan menulis
$g_\varepsilon$ lewat pembalikan Fourier transformasinya
(\cref{exo:b3:fouriertransform:4}, yang diskalakan ulang):
$g_\varepsilon(u) = \frac1{2\pi}\int
\eu^{-\varepsilon^2\xi^2/2}\eu^{\iu\xi u}\dd\xi$, lalu Fubini
(sebab semuanya terdominasi faktor Gaussnya):
\[
p_\varepsilon(x) = \frac1{2\pi}\int_\R
\varphi_X(-\xi)\,\eu^{-\varepsilon^2\xi^2/2}\,
\eu^{\iu\xi x}\,\dd\xi ,
\]
yakni fungsional $\varphi_X$ semata. Jika $\varphi_X =
\varphi_Y$: maka $X + \varepsilon N$ dan $Y + \varepsilon N$
berdistribusi sama bagi setiap $\varepsilon$; sedangkan bagi $f$ yang
\omterm{def:b3:topology:continuity}{kontinu} terbatas, $\E f(X + \varepsilon N) \to \E f(X)$ saat $\varepsilon
\to 0$ (menurut konvergensi terdominasi, sebab $X + \varepsilon N \to X$
titik demi titik pada ruang hasil kalinya), sehingga $\E f(X) = \E f(Y)$ bagi
setiap $f$ semacam itu --- dan ini menentukan distribusinya: sebab bagi setiap $t$,
apitlah $\mathbf 1_{\intoc{-\infty}t}$ di antara landaian
\omterm{def:b3:topology:continuity}{kontinu} terbatas $f_k^\pm$ (yang sama dengan $1$ pada
$\intoc{-\infty}{t \mp \frac1k}$, sama dengan $0$ di luar $t \pm
\frac1k$, dan afin di antaranya); lalu melimitkan $\E
f_k^-(X) \leq F_X(t) \leq \E f_k^+(X)$ memberikan $F_X(t) =
F_Y(t)$ pada setiap $t$ yang keduanya \omterm{def:b3:topology:continuity}{kontinu} di sana, sehingga
di mana-mana menurut \omterm{def:b3:topology:continuity}{kekontinuan} kanan dan \omterm{ex:b3:lebesgue:gamma}{kepadatan} titik
\omterm{def:b3:topology:continuity}{kekontinuan} bersamanya (sebab kedua $F$ berlompatan terbilang banyak);
dan fungsi distribusi yang sama memaksa distribusi yang sama
(\cref{exo:b3:measure:3}, yang bersandar pada
\cref{thm:b3:measure:uniqueness}).
\end{proof}

\section{Konvergensi dalam distribusi}

\begin{definition}\label{def:b3:clt:cid}
Variabel $X_n$ \emph{konvergen dalam distribusi} (atau dalam hukumnya) ke $X$,
yang ditulis $X_n \Rightarrow X$, jika\index{konvergensi dalam
distribusi}
\[
\E\bigl[f(X_n)\bigr] \longrightarrow \E\bigl[f(X)\bigr]
\qquad\text{bagi setiap fungsi kontinu terbatas } f\colon\R\to\R .
\]
Secara setara (\cref{exo:b3:clt:4}): $F_{X_n}(t) \to F_X(t)$
pada setiap titik \omterm{def:b3:topology:continuity}{kekontinuan} $t$ milik $F_X$. Lalu $X_n$ tak perlu
hidup pada \omterm{def:b3:probability:space}{ruang peluang} bersama: sebab hanya distribusinya yang penting.
\end{definition}

\begin{theorem}[Teorema pemilihan Helly]
\label{thm:b3:clt:helly}
Setiap barisan $(F_n)$ fungsi distribusi mempunyai
subbarisan yang konvergen titik demi titik, pada setiap titik \omterm{def:b3:topology:continuity}{kekontinuan}
limitnya, ke $G \colon \R \to \intcc01$ yang tak turun dan \omterm{def:b3:topology:continuity}{kontinu}
kanan --- mungkin dengan $G(+\infty) - G(-\infty) <
1$ (sebab massanya boleh lolos ke tak hingga).\index{teorema Helly}
\end{theorem}

\begin{proof}
Ekstraksi diagonalnya memberikan $F_{n_k}(q) \to \ell(q)$ bagi setiap
$q$ rasional (sebab nilainya di $\intcc01$ yang \omterm{def:b3:topology:compact}{kompak}). Definisikanlah
$G(t) = \inf\{\ell(q) : q \in \Q, q > t\}$: yang tak turun;
dan \omterm{def:b3:topology:continuity}{kontinu} kanan (sebab infimum atas \omterm{def:b3:topology:topology}{persekitaran} rasional yang
mengecil dari kanan). Lalu pada titik \omterm{def:b3:topology:continuity}{kekontinuan} $t$ milik
$G$: bagi bilangan rasional $q_1 < t < q_2$,
\[
\ell(q_1) \leq \liminf F_{n_k}(t) \leq \limsup F_{n_k}(t)
\leq \ell(q_2),
\]
menurut kemonotonan setiap $F_{n_k}$. Lalu dari definisi
$G$ sebagai infimum dan kemonotonan $\ell$ pada
bilangan rasionalnya: $G(s) \leq \ell(q) \leq G(q)$ setiap kali $s < q$.
Lalu mengambil $s < q_1 < t$ memberikan $\ell(q_1) \geq G(s)$, dan
$\ell(q_2) \leq G(q_2)$; sehingga dengan membiarkan $s \uparrow t$ dan $q_2
\downarrow t$, \omterm{def:b3:topology:continuity}{kekontinuan} $G$ di $t$ mengapit baik
$\liminf$ maupun $\limsup$-nya ke $G(t)$.
\end{proof}

\begin{lemma}[Keketatan dari fungsi karakteristiknya]
\label{lem:b3:clt:tightness}
Bagi sembarang \omterm{def:b3:probability:space}{peubah acak} $X$ dan $u > 0$:
\[
\P\Bigl(\abs X \geq \frac2u\Bigr) \;\leq\;
\frac1u\int_{-u}^{u}\bigl(1 -
\operatorname{Re}\varphi_X(\xi)\bigr)\,\dd\xi .
\]
\end{lemma}

\begin{proof}
Menurut Tonelli--Fubini (sebab integrannya terbatas dan daerahnya berhingga terhadap
$\xi$):
\[
\frac1u\int_{-u}^u\bigl(1 -
\operatorname{Re}\varphi_X(\xi)\bigr)\dd\xi
= \E\Bigl[\frac1u\int_{-u}^u(1 - \cos(\xi X))\,\dd\xi\Bigr]
= 2\,\E\Bigl[1 - \frac{\sin(uX)}{uX}\Bigr]
\]
(tafsirkanlah kurungnya sebagai limitnya $0$ di $X = 0$). Sedangkan
integrannya tak negatif (sebab $\abs{\sin t} \leq \abs t$), dan
bagi $\abs{uX} \geq 2$: $1 - \frac{\sin(uX)}{uX} \geq 1 -
\frac1{\abs{uX}} \geq \frac12$. Jadi menyimpan hanya kejadian
$\{\abs{uX} \geq 2\}$ di dalam \omterm{def:b3:probability:space}{nilai harapannya}
menyisakan setidaknya $2 \cdot \frac12\,\P(\abs X \geq \frac2u)$,
dan itulah klaimnya.
\end{proof}

\begin{theorem}[Teorema kekontinuan Lévy]
\label{thm:b3:clt:levy}
Misalkan $(X_n)$ \omterm{def:b3:probability:space}{peubah acak} yang \omterm{def:b3:clt:cf}{fungsi
karakteristiknya} konvergen titik demi titik: $\varphi_{X_n}(\xi) \to
\varphi(\xi)$ bagi setiap $\xi$, dengan $\varphi =
\varphi_X$ sebagai \omterm{def:b3:clt:cf}{fungsi karakteristik} suatu \omterm{def:b3:probability:space}{peubah
acak} $X$. Maka $X_n \Rightarrow X$.\index{teorema
kekontinuan Lévy}
\end{theorem}

\begin{proof}
\emph{Keketatannya.} Tetapkanlah $\varepsilon > 0$. Karena $\varphi$
\omterm{def:b3:topology:continuity}{kontinu} di $0$ dengan $\varphi(0) = 1$, pilihlah $u > 0$ dengan
$\frac1u\int_{-u}^u(1 - \operatorname{Re}\varphi) <
\varepsilon$; lalu menurut konvergensi terdominasi (sebab integrannya terbatas oleh
$2$ pada $[-u,u]$ yang tetap), integral yang sama bagi
$\varphi_{X_n}$ bernilai $< 2\varepsilon$ bagi $n$ yang besar:
sehingga \cref{lem:b3:clt:tightness} memberikan $\P(\abs{X_n} \geq \frac2u)
\leq 2\varepsilon$ bagi $n$ yang besar, sedangkan membesarkan konstantanya
menangani berhingga banyak sisanya: jadi distribusinya \emph{ketat}
--- tak ada massa yang lolos.

\emph{Subbarisannya.} Misalkan $(F_{n_k})$ sembarang subbarisan; lalu menurut
Helly (\cref{thm:b3:clt:helly}) sarikanlah $F_{n_{k_j}} \to G$
pada titik \omterm{def:b3:topology:continuity}{kekontinuannya}. Keketatannya memaksa $G(-\infty) = 0$,
$G(+\infty) = 1$ (sebab $G(\frac2u) - G(-\frac2u) \geq 1 -
2\varepsilon$ pada titik \omterm{def:b3:topology:continuity}{kekontinuannya}): jadi $G$ merupakan fungsi
distribusi yang sejati, milik suatu \omterm{def:b3:probability:space}{peubah acak} $Y$. Maka
$X_{n_{k_j}} \Rightarrow Y$ (\cref{exo:b3:clt:4},
yakni konvergensi distribusinya dari $F$-nya), sehingga
$\varphi_{X_{n_{k_j}}} \to \varphi_Y$ \emph{titik demi titik} (sebab $x
\mapsto \eu^{\iu\xi x}$ \omterm{def:b3:topology:continuity}{kontinu} terbatas, pada bagian real dan
imajinernya secara terpisah); lalu dengan membandingkannya dengan hipotesisnya:
$\varphi_Y = \varphi = \varphi_X$, sehingga keinjektifannya
(\cref{thm:b3:clt:injectivity}) memberikan $Y \sim X$, yakni $G =
F_X$.

\emph{Kesimpulannya.} Setiap subbarisan $(F_n)$ mempunyai
subsubbarisan yang konvergen ke $F_X$ yang \emph{sama} (pada titik
\omterm{def:b3:topology:continuity}{kekontinuannya}); sehingga $F_n(t) \to F_X(t)$ pada setiap
titik \omterm{def:b3:topology:continuity}{kekontinuan} $t$ (sebab barisan real yang setiap
subbarisannya mempunyai subsubbarisan berlimit sama akan
konvergen): jadi $X_n \Rightarrow X$.
\end{proof}

\section{Teorema limit pusat}

\begin{theorem}[Teorema limit pusat]\label{thm:b3:clt:clt}
Misalkan $(X_n)$ i.i.d.\ dengan $\E X_1 = m$ dan $\V(X_1) =
\sigma^2 \in \intoo0\infty$. Maka\index{teorema limit
pusat}
\[
\frac{S_n - nm}{\sigma\sqrt n} \;\Longrightarrow\; \mathcal
N(0, 1) :
\qquad
\P\Bigl(a \leq \frac{S_n - nm}{\sigma\sqrt n} \leq
b\Bigr) \longrightarrow
\frac{1}{\sqrt{2\pi}}\int_a^b\eu^{-x^2/2}\,\dd x
\]
bagi setiap $a < b$.
\end{theorem}

\begin{proof}
Pusatkan lalu normalkan: $Z_i = \frac{X_i - m}{\sigma}$
(yang i.i.d., berrata-rata $0$ dan bervarians $1$) dan $T_n =
\frac1{\sqrt n}\sum_{i\leq n}Z_i$. Lalu menurut \omterm{def:b3:probability:independence}{kebebasannya} dan
aturan afinnya (\cref{prop:b3:clt:cfbasics}):
\[
\varphi_{T_n}(\xi) =
\varphi_{Z}\Bigl(\frac{\xi}{\sqrt n}\Bigr)^{n},
\qquad
\varphi_Z(\eta) = 1 - \frac{\eta^2}2 + \eta^2\rho(\eta),\quad
\rho(\eta)\to0 .
\]
Tetapkanlah $\xi$ lalu ambil $a_n = \varphi_Z(\xi/\sqrt n)$, $b_n = 1 -
\frac{\xi^2}{2n}$: keduanya bermodulus $\leq 1$ bagi $n$ yang besar
(sebab $\abs{b_n} \leq 1$ begitu $\xi^2 \leq 4n$; sedangkan $\abs{a_n} \leq 1$
selalu). Lalu ketaksamaan elementernya $\abs{a^n - b^n} \leq
n\abs{a - b}$ bagi $\abs a, \abs b \leq 1$ (lewat teleskop
$a^n - b^n = \sum a^k(a - b)b^{n-1-k}$) memberikan
\[
\Bigl|\varphi_{T_n}(\xi) - \Bigl(1 -
\frac{\xi^2}{2n}\Bigr)^{n}\Bigr|
\leq n\,\Bigl|\varphi_Z\Bigl(\frac\xi{\sqrt n}\Bigr) - 1 +
\frac{\xi^2}{2n}\Bigr|
= \xi^2\,\Bigl|\rho\Bigl(\frac{\xi}{\sqrt n}\Bigr)\Bigr|
\longrightarrow 0,
\]
sedangkan $\bigl(1 - \frac{\xi^2}{2n}\bigr)^n \to
\eu^{-\xi^2/2}$ (lewat logaritma realnya). Jadi $\varphi_{T_n}(\xi) \to
\eu^{-\xi^2/2} = \varphi_{\mathcal N(0,1)}(\xi)$
(\cref{prop:b3:clt:cfbasics}(d)) bagi setiap $\xi$: sehingga Lévy
(\cref{thm:b3:clt:levy}) menyimpulkan $T_n \Rightarrow \mathcal
N(0,1)$. Adapun peluang selangnya menyusul sebab $F_{\mathcal
N}$ \omterm{def:b3:topology:continuity}{kontinu} di mana-mana.
\end{proof}

\begin{example}[Selang kepercayaan, yang diturunkan secara jujur]
\label{ex:b3:clt:confidence}
Jajakilah $n$ pemilih yang bebas; maka $\hat p_n = S_n/n$ menaksir
$p$ yang sebenarnya, dengan $\sigma^2 = p(1-p) \leq \frac14$. Lalu teorema limit
pusatnya memberikan, bagi $n$ yang besar,
\[
\P\Bigl(\abs{\hat p_n - p} \leq
\frac{z}{2\sqrt n}\Bigr)
\;\geq\; \P\Bigl(\Bigl|\frac{S_n - np}{\sigma\sqrt n}\Bigr|
\leq z\Bigr)
\longrightarrow \Phi(z) - \Phi(-z),
\]
dengan $\Phi$ sebagai fungsi \omterm{ex:b3:probability:laws}{distribusi Gauss} bakunya.
Dengan $z = 1.96$: kepercayaan asimtotiknya $95\%$, sedangkan marjin
$\frac{1.96}{2\sqrt n} \leq 3\%$ menuntut $n \geq
\bigl(\frac{1.96}{0.06}\bigr)^2 \approx 1068$ --- yakni bilangan
di balik setiap ``$\pm3$ poin, $95\%$'' yang orang baca; bandingkanlah
dengan $5556$ milik Chebyshev (\cref{exo:b3:probability:7}). Adapun
$\sqrt n$-nya universal: bahwa untuk memarukan galatnya, cuplikannya harus
dilipatempatkan --- yakni hukum yang sama yang menetapkan ongkos \omterm{ex:b3:probability:sllnapps}{Monte Carlo}
(\cref{exo:b3:clt:7}).
\end{example}

\section{Vektor Gauss}

\begin{definition}\label{def:b3:clt:gaussianvector}
Sebuah vektor acak $X = (X_1, \dots, X_d)$ disebut
\emph{Gauss}\index{vektor Gauss} jika setiap kombinasi
linearnya $\langle t, X\rangle = \sum t_iX_i$ merupakan
variabel Gauss real (yang mungkin merosot). Distribusinya
ditentukan oleh vektor rata-ratanya $m = (\E X_i)$ dan
\emph{matriks kovarians} $\Sigma = \bigl(\operatorname{Cov}
(X_i, X_j)\bigr)$: sebab \omterm{def:b3:clt:cf}{fungsi karakteristik}
vektornya, $\varphi_X(t) = \E\eu^{\iu\langle t, X\rangle}$, adalah
nilai di $1$ bagi \omterm{def:b3:clt:cf}{fungsi karakteristik} $\langle t, X\rangle$:
\[
\varphi_X(t) = \exp\Bigl(\iu\langle t, m\rangle -
\tfrac12\,t^{\mathsf T}\Sigma\,t\Bigr),
\]
sedangkan \omterm{def:b3:clt:cf}{fungsi karakteristik} berdimensi $d$ bersifat injektif
(lewat bukti pemulusan yang sama dengan \cref{thm:b3:clt:injectivity},
dengan Gauss koordinat demi koordinat).
\end{definition}

\begin{theorem}\label{thm:b3:clt:gaussianvector}
Misalkan $X$ sebuah \omterm{def:b3:clt:gaussianvector}{vektor Gauss}.
\begin{enumerate}
  \item Setiap peta afin $AX + b$ merupakan \omterm{def:b3:clt:gaussianvector}{vektor Gauss}.
  \item Komponen $X_i$ bersifat \emph{bebas} jika dan
        hanya jika $\Sigma$ diagonal: jadi bagi variabel \omterm{def:b3:clt:gaussianvector}{Gauss}
        bersama, tak berkorelasi $=$ bebas.
  \item Jika $\Sigma$ terbalikkan, maka $X$ berkepadatan
        $\frac{1}{(2\pi)^{d/2}\sqrt{\det\Sigma}}
        \exp\bigl(-\frac12(x - m)^{\mathsf T}\Sigma^{-1}(x -
        m)\bigr)$.
\end{enumerate}
\end{theorem}

\begin{proof}
(1) Kombinasi linear komponen $AX + b$ merupakan fungsi afin
kombinasi linear $X$: jadi \omterm{def:b3:clt:gaussianvector}{Gauss} (sebab peta afin
sebuah variabel \omterm{def:b3:clt:gaussianvector}{Gauss} bersifat \omterm{def:b3:clt:gaussianvector}{Gauss}).
(2) Jika $\Sigma$ diagonal, maka \omterm{def:b3:clt:cf}{fungsi karakteristiknya}
memfaktor: $\varphi_X(t) = \prod_i\exp(\iu t_im_i -
\frac12\Sigma_{ii}t_i^2) = \prod\varphi_{X_i}(t_i)$, yang merupakan
\omterm{def:b3:clt:cf}{fungsi karakteristik} distribusi hasil kalinya
(\cref{thm:b3:probability:independence} yang dibaca lewat
keinjektifan berdimensi $d$): sehingga komponennya bebas.
Sedangkan konversnya adalah lenyapnya kovarians variabel $L^2$
yang bebas.
(3) Diagonalkanlah $\Sigma = P D P^{\mathsf T}$ (dengan $P$ ortogonal dan
$D > 0$ diagonal --- \cref{exo:b3:submanifolds:8}); maka
vektor $Y = P^{\mathsf T}(X - m)$ bersifat \omterm{def:b3:clt:gaussianvector}{Gauss} dengan
kovarians $D$: sehingga menurut (2) komponennya bebas dan
$\mathcal N(0, d_i)$, jadi $Y$ berkepadatan hasil kali; lalu dorong
majulah lewat $x = m + PY$ yang mengawetkan volume
(\cref{thm:b3:product:linearchange}, sebab $\abs{\det P} = 1$) lalu
tulislah kembali eksponennya secara invarian.
\end{proof}

\begin{theorem}[Teorema limit pusat berdimensi banyak]
\label{thm:b3:clt:multiclt}
Misalkan $(X_n)$ merupakan \emph{vektor} acak i.i.d.\ yang
\omterm{def:b3:lebesgue:l1}{terintegralkan} kuadrat di $\R^d$ dengan rata-rata $m$ dan matriks kovarians
$\Sigma$. Maka $\frac{S_n - nm}{\sqrt n}$ konvergen dalam
distribusi ke \omterm{def:b3:clt:gaussianvector}{vektor Gauss} $\mathcal N(0,
\Sigma)$.\index{teorema limit pusat!berdimensi banyak}
\admitted
\end{theorem}

\begin{remark}\label{rem:b3:clt:multiclt}
Hampir semuanya sudah ada di tangan kita. Sebab bagi setiap
arah $t \in \R^d$, variabel real $\langle t,
\frac{S_n - nm}{\sqrt n}\rangle$ merupakan jumlah ternormalkan
variabel real i.i.d.\ bervarians $t^{\mathsf T}\Sigma t$,
sehingga perhitungan \cref{thm:b3:clt:clt} memberikan konvergensi
titik demi titik \omterm{def:b3:clt:cf}{fungsi karakteristik} berdimensi $d$-nya
ke $\eu^{-t^{\mathsf T}\Sigma t/2}$, yakni \omterm{def:b3:clt:cf}{fungsi
karakteristik} $\mathcal N(0, \Sigma)$
(\cref{def:b3:clt:gaussianvector}). Adapun yang belum kita
buktikan ulang adalah teorema \omterm{def:b3:topology:continuity}{kekontinuan} Lévy \emph{di $\R^d$}:
sebab pemilihan Helly dan taksiran keketatannya diperumum
secara rutin (koordinat demi koordinat), dan reduksi Cramér--Wold
ini dikerjakan dengan jujur pada setiap kuliah peluang
pascasarjana; jadi tak ada yang diperlukan di luar metode bab ini.
\end{remark}

\begin{method}\label{met:b3:clt:method}
Untuk mengenali sebuah distribusi limit: hitunglah \omterm{def:b3:clt:cf}{fungsi karakteristiknya},
ambillah limit titik demi titiknya, kenalilah (\omterm{def:b3:clt:gaussianvector}{Gauss}
$\eu^{-\sigma^2\xi^2/2}$, Poisson
$\eu^{\lambda(\eu^{\iu\xi}-1)}$, eksponensial $\frac{\lambda}
{\lambda - \iu\xi}$, \dots) lalu panggillah Lévy. Ritual tiga
langkahnya (\omterm{def:b3:probability:independence}{kebebasan} $\to$ hasil kali; Taylor di $0$ $\to$
limit eksponensial; Lévy $\to$ konvergensi distribusi) membuktikan
teorema limit pusatnya, hukum kejadian langka Poisson
(\cref{exo:b3:clt:5}), dan setiap teorema limit klasik pada
kuliah ini. Sedangkan untuk pernyataan h.p., kembalilah ke
perkakas \cref{ch:b3:probability}: sebab kedua babnya menjawab
pertanyaan berbeda tentang $S_n$ yang sama.
\end{method}

\section{Latihan}

\begin{exercise}[$\star$]\label{exo:b3:clt:1}
Hitunglah \omterm{def:b3:clt:cf}{fungsi karakteristiknya}: seragam pada
$\intcc{-1}1$; eksponensial $\mathcal E(\lambda)$; Poisson
$\mathcal P(\lambda)$; binomial $\mathcal B(n, p)$. Lalu simpulkanlah
lewat \cref{thm:b3:clt:injectivity} bahwa jumlah variabel Poisson
yang bebas (dengan $\lambda, \mu$) bersifat Poisson $(\lambda +
\mu)$.
\end{exercise}

\begin{exercise}[$\star\star$]\label{exo:b3:clt:2}
(a) Tunjukkanlah bahwa $\varphi_X$ bernilai real jika dan hanya jika $X$
dan $-X$ berdistribusi sama (yakni variabel yang \emph{setangkup}).
(b) Andaikanlah $\abs{\varphi_X(\xi_0)} = 1$ bagi suatu $\xi_0 \neq
0$. Tunjukkanlah bahwa $X$ hampir pasti bertumpu pada barisan
aritmetika $a + \frac{2\pi}{\xi_0}\Z$ \emph{(tulislah
$\varphi_X(\xi_0) = \eu^{\iu\theta}$ lalu hitung $\E[1 -
\cos(\xi_0X - \theta)]$)}. Lalu simpulkanlah bahwa jika $X$ berkepadatan,
maka $\abs{\varphi_X(\xi)} < 1$ bagi setiap $\xi \neq 0$.
\end{exercise}

\begin{exercise}[$\star\star$]\label{exo:b3:clt:3}
Misalkan $X \sim \mathcal N(m_1, \sigma_1^2)$ dan $Y \sim
\mathcal N(m_2, \sigma_2^2)$ saling bebas. Tunjukkanlah $X + Y
\sim \mathcal N(m_1 + m_2, \sigma_1^2 + \sigma_2^2)$, dan
lebih umum bahwa keluarga Gaussnya stabil terhadap
jumlah yang bebas dan pemetaan afin. Pertentangkanlah: apakah jumlah
dua variabel \omterm{def:b3:clt:gaussianvector}{Gauss} yang \emph{bergantung} selalu \omterm{def:b3:clt:gaussianvector}{Gauss}?
(\cref{exo:b3:clt:9}.)
\end{exercise}

\begin{exercise}[$\star\star$]\label{exo:b3:clt:4}
(a) Buktikanlah kesetaraan pada \cref{def:b3:clt:cid}: bahwa jika
$\E f(X_n) \to \E f(X)$ bagi setiap $f$ yang \omterm{def:b3:topology:continuity}{kontinu} terbatas,
maka $F_{X_n}(t) \to F_X(t)$ pada titik \omterm{def:b3:topology:continuity}{kekontinuannya}
\emph{(apitlah $\mathbf 1_{\intoc{-\infty}t}$ di antara dua
landaian-tangga yang \omterm{def:b3:topology:continuity}{kontinu})}; dan sebaliknya
\emph{(hampirilah $f$ yang \omterm{def:b3:topology:continuity}{kontinu} terbatas oleh jumlah fungsi
landaian, atau syaratkanlah pada kisi halus titik
\omterm{def:b3:topology:continuity}{kekontinuannya})} --- dan konversnya boleh ditangani bagi $f$ yang \omterm{def:b3:topology:continuity}{kontinu}
seragam lebih dahulu, lalu secara umum.
(b) Tunjukkanlah bahwa $X_n \Rightarrow c$ (yakni sebuah konstanta) mengakibatkan $X_n
\to c$ dalam peluang.
\end{exercise}

\begin{exercise}[$\star\star$]\label{exo:b3:clt:5}
(Hukum kejadian langka) Misalkan $X_n \sim \mathcal B(n, p_n)$ dengan
$np_n \to \lambda > 0$. Tunjukkanlah, lewat \omterm{def:b3:clt:cf}{fungsi karakteristiknya}
dan \cref{thm:b3:clt:levy}, bahwa $X_n \Rightarrow \mathcal
P(\lambda)$. Lalu sebagai uji kewarasan numerik: bandingkanlah $\P(X = 0)$ bagi
$\mathcal B(100, 0.02)$ dan $\mathcal P(2)$.
\end{exercise}

\begin{exercise}[$\star\star$]\label{exo:b3:clt:6}
(a) Sebuah dadu adil dilempar $n = 1000$ kali; hampirilah
peluang bahwa totalnya melampaui $3600$ (dengan rata-rata $3500$ dan
varians per lemparan $\frac{35}{12}$).
(b) Untuk $S \sim \mathcal B(100, \frac12)$, hampirilah
$\P(45 \leq S \leq 55)$ lewat teorema limit pusatnya dengan koreksi
\omterm{def:b3:topology:continuity}{kekontinuan} ($\pm\frac12$), lalu berilah komentar atas pengaruh
koreksinya.
\end{exercise}

\begin{exercise}[$\star\star$]\label{exo:b3:clt:7}
(Galat \omterm{ex:b3:probability:sllnapps}{Monte Carlo}) Pada latar
\cref{pb:b3:probability:1}, pertanyaan 11, dengan $g \in
L^2(\intcc01^d)$, ambillah $\sigma^2 = \V(g(U_1))$ dan $I = \int
g$. Tunjukkanlah
\[
\sqrt n\,\Bigl(\frac1n\sum_{k\leq n}g(U_k) - I\Bigr)
\Longrightarrow \mathcal N(0, \sigma^2),
\]
lalu simpulkanlah batang galat $95\%$ asimtotiknya $\pm
1.96\,\sigma/\sqrt n$ --- yang tak bergantung pada dimensi $d$.
Bandingkanlah dengan kaidah titik tengah deterministik pada dimensi
$d$ (yang galatnya $\sim n^{-2/d}$ bagi integran $\mathcal C^2$):
mulai dimensi berapa pencuplikan acaknya menang?
\end{exercise}

\begin{exercise}[$\star\star\star$]\label{exo:b3:clt:8}
(Slutsky) Andaikanlah $X_n \Rightarrow X$ dan $Y_n \to c$ dalam
peluang (dengan $c$ konstanta). Tunjukkanlah $X_n + Y_n \Rightarrow X +
c$ dan $Y_nX_n \Rightarrow cX$. \emph{(Bekerjalah dengan
\omterm{def:b3:clt:cf}{fungsi karakteristiknya} dan batas $\abs{\E\eu^{\iu\xi
(X_n+Y_n)} - \eu^{\iu\xi c}\E\eu^{\iu\xi X_n}} \leq
\E\abs{\eu^{\iu\xi(Y_n - c)} - 1}$, lalu belahlah pada $\abs{Y_n - c}
\leq \delta$.)} Penerapannya: pada
\cref{ex:b3:clt:confidence}, benarkanlah penggantian
$\sigma = \sqrt{p(1-p)}$ yang tak diketahui oleh $\sqrt{\hat p_n(1 - \hat
p_n)}$.
\end{exercise}

\begin{exercise}[$\star\star\star$]\label{exo:b3:clt:9}
Misalkan $X \sim \mathcal N(0,1)$ dan $\varepsilon$ bebas
dengan $\P(\varepsilon = \pm1) = \frac12$; tetapkanlah $Y =
\varepsilon X$.
(a) Tunjukkanlah $Y \sim \mathcal N(0,1)$ dan
$\operatorname{Cov}(X, Y) = 0$.
(b) Tunjukkanlah bahwa $X$ dan $Y$ \emph{tidak} bebas, dan bahwa
$(X, Y)$ bukan \omterm{def:b3:clt:gaussianvector}{vektor Gauss} \emph{(hitunglah $\P(X + Y =
0)$)}.
(c) Pelajarannya: \cref{thm:b3:clt:gaussianvector}(2) menuntut
sifat \omterm{def:b3:clt:gaussianvector}{Gauss} bersamanya --- sebab ``\omterm{def:b3:clt:gaussianvector}{Gauss} yang tak berkorelasi'' saja
tak membuktikan apa pun.
\end{exercise}

\begin{exercise}[$\star\star$]\label{exo:b3:clt:10}
Distribusi Cauchy berkepadatan $\frac1{\pi(1 + x^2)}$.
(a) Tunjukkanlah bahwa \omterm{def:b3:clt:cf}{fungsi karakteristiknya} $\eu^{-\abs\xi}$
(\cref{exo:b3:fouriertransform:1} dan pembalikannya).
(b) Tunjukkanlah bahwa jika $X_1, \dots, X_n$ i.i.d.\ Cauchy, maka
$\frac{S_n}n$ kembali Cauchy --- yakni distribusi yang \emph{sama}: sehingga
rata-ratanya tak pernah memusat.
(c) Damaikanlah dengan hukum bilangan besar dan teorema limit pusatnya:
hipotesis mana yang gagal? (Hitunglah $\E\abs{X_1}$.)
\end{exercise}

\begin{exercise}[$\star\star$]\label{exo:b3:clt:11}
(Distribusi stabil dalam kandungan) Misalkan $(X_n)$ i.i.d.\ Cauchy
baku (\cref{exo:b3:clt:10}).
(a) Tunjukkanlah bahwa bagi sembarang $a, b > 0$, $aX_1 + bX_2$ berdistribusi
seperti $(a + b)X_1$: jadi keluarga Cauchy bersifat \emph{stabil
sejati} berindeks $1$.
(b) Tunjukkanlah bahwa keluarga Gaussnya stabil sejati berindeks
$2$: yakni $aX_1 + bX_2 \sim \sqrt{a^2 + b^2}\,X_1$ bagi
$X_i$ yang i.i.d.\ $\mathcal N(0,1)$.
(c) Jelaskanlah, lewat \omterm{def:b3:clt:cf}{fungsi karakteristik} berbentuk
$\eu^{-c\abs\xi^\alpha}$, mengapa kestabilan berindeks-$\alpha$
memaksa penormalan $n^{1/\alpha}$ bagi jumlahnya, dan apa
artinya bagi cekungan tarikan teorema limit pusatnya: yakni jumlah
i.i.d.\ mana yang dapat konvergen, setelah penormalan afin, ke
distribusi Cauchy alih-alih ke \omterm{def:b3:clt:gaussianvector}{Gauss}?
\end{exercise}

\begin{exercise}[$\star\star$]\label{exo:b3:clt:12}
(Fungsi distribusi empiris) Misalkan $(X_n)$ i.i.d.\
berfungsi distribusi $F$, dan $F_n(t) =
\frac1n\#\{k \leq n : X_k \leq t\}$.
(a) Tetapkanlah $t$. Tunjukkanlah bahwa $n F_n(t) \sim \mathcal B(n, F(t))$,
bahwa $F_n(t) \to F(t)$ secara h.p.\ (\cref{thm:b3:probability:slln}),
dan bahwa
\[
\sqrt n\,\bigl(F_n(t) - F(t)\bigr) \Longrightarrow
\mathcal N\bigl(0,\ F(t)(1 - F(t))\bigr) .
\]
(b) Pada $t$ manakah varians asimtotiknya maksimal?
Tafsirkanlah: bahwa medianlah tempat sebuah distribusi empiris
paling sukar dipastikan.
(c) Untuk $F$ yang \omterm{def:b3:topology:continuity}{kontinu}, tunjukkanlah bahwa distribusi
$\sup_t\abs{F_n(t) - F(t)}$ tak bergantung pada $F$
\emph{(reduksikanlah ke variabel seragam lewat
\cref{exo:b3:probability:1})} --- yakni mukjizat bebas-distribusi
di balik uji Kolmogorov--Smirnov; dan tak diminta perhitungan
distribusi itu.
\end{exercise}

\section{Soal: bukti Lindeberg bagi teorema limit pusatnya, dengan
laju}

\begin{problem}[{Soal akhir pekan --- metode penggantian}]
\label{pb:b3:clt:1}
Lindeberg (1922) membuktikan teorema limit pusatnya lewat gagasan
yang sederhananya melucuti: \emph{tukarlah sukunya satu per satu
dengan variabel \omterm{def:b3:clt:gaussianvector}{Gauss}} lalu kendalikan setiap penukaran lewat uraian
Taylor. Metodenya tak memerlukan analisis Fourier, menghasilkan
laju galat yang eksplisit, dan hari ini menjalankan bukti keuniversalan
di seluruh teori peluang. Misalkan $(X_i)$ i.i.d., terpusat, dengan
$\V(X_1) = 1$ dan $\beta = \E\abs{X_1}^3 < \infty$; lalu misalkan
$(N_i)$ i.i.d.\ $\mathcal N(0,1)$, yang bebas dari
$X_i$-nya (keberadaannya: \cref{thm:b3:probability:existence}). Tetapkanlah
\[
T_n = \frac{X_1 + \dots + X_n}{\sqrt n},
\qquad
G_n = \frac{N_1 + \dots + N_n}{\sqrt n} \sim \mathcal N(0,1).
\]

\medskip
\noindent\textbf{Bagian I --- Identitas penukarannya.} Tetapkanlah $f
\in \mathcal C^3_b(\R)$ (yakni tiga turunan \omterm{def:b3:topology:continuity}{kontinu} yang
terbatas; dengan $M_3 = \sup\abs{f'''}$). Lalu bagi $0 \leq i \leq n$
definisikanlah jumlah hibridanya
\[
H_i = \frac{X_1 + \dots + X_i + N_{i+1} + \dots +
N_n}{\sqrt n},
\]
sehingga $H_n = T_n$ dan $H_0 = G_n$.
\begin{enumerate}
  \item Tulislah $H_i = W_i + \frac{X_i}{\sqrt n}$ dan $H_{i-1}
        = W_i + \frac{N_i}{\sqrt n}$ dengan $W_i =
        \frac{1}{\sqrt n}\bigl(\sum_{j<i}X_j +
        \sum_{j>i}N_j\bigr)$, lalu perhatikan bahwa $W_i$
        bebas dari pasangan $(X_i, N_i)$. Benarkanlah.
  \item Taylor dengan sisa integral atau sisa Lagrange: bagi sembarang
        $w, h$ yang real:
        \[
        \Bigl|f(w + h) - f(w) - f'(w)h -
        \tfrac12f''(w)h^2\Bigr| \leq
        \frac{M_3\,\abs h^3}{6} .
        \]
  \item Terapkanlah pertanyaan 2 dua kali (dengan $h = \frac{X_i}{\sqrt n}$ dan
        $h = \frac{N_i}{\sqrt n}$ di $w = W_i$), ambillah
        \omterm{def:b3:probability:space}{nilai harapannya}, lalu pakailah \omterm{def:b3:probability:independence}{kebebasannya} ditambah kecocokan
        dua momen pertama $X_i$ dan $N_i$ untuk menunjukkan
        \[
        \bigl|\E f(H_i) - \E f(H_{i-1})\bigr|
        \leq \frac{M_3}{6}\cdot
        \frac{\beta + \gamma}{n^{3/2}},
        \qquad \gamma = \E\abs{N_1}^3 =
        \frac{2\sqrt2}{\sqrt\pi} .
        \]
  \item Teleskopkanlah atas $i$ lalu simpulkan \emph{batas
        Lindeberg}-nya:
        \[
        \bigl|\E f(T_n) - \E f(G_n)\bigr| \leq
        \frac{M_3\,(\beta + \gamma)}{6\,\sqrt n} .
        \]
\end{enumerate}

\medskip
\noindent\textbf{Bagian II --- Dari $f$ yang mulus ke teorema limit pusatnya.}
\begin{enumerate}[resume]
  \item Tunjukkanlah bahwa $\E f(T_n) \to \E f(N)$ bagi setiap $f \in
        \mathcal C_b^3$, lalu naikkanlah ke setiap $f$ yang \omterm{def:b3:topology:continuity}{kontinu}
        terbatas: yakni diberikan $f$ semacam itu dan $\varepsilon$,
        bangunlah $f_\varepsilon \in \mathcal C^3_b$ dengan
        $\norm{f - f_\varepsilon}_\infty \leq \varepsilon$
        pada selang yang besar --- misalnya konvolusikanlah $f$ dengan
        gundukan $\mathcal C^\infty$
        (\cref{thm:b3:lp:regularization}) --- lalu tanganilah
        ekornya lewat keketatannya ($\V(T_n) = 1$ dan Chebyshev).
        Simpulkanlah $T_n \Rightarrow \mathcal N(0, 1)$: yakni
        teorema limit pusatnya, yang dibuktikan ulang.
  \item Di manakah buktinya memakai bahwa $X_i$ berdistribusi
        \emph{identik}? Tunjukkanlah bahwa ia nyaris tak
        memakainya: nyatakan lalu buktikanlah versinya bagi $X_i$ yang
        bebas, terpusat, dan tak identik dengan
        $\sum_i\V(X_i) = s_n^2$ beserta momen ketiganya, lalu perolehlah
        galat $\frac{M_3}{6s_n^3}\sum_i\bigl(\E\abs{X_i}^3
        + \V(X_i)^{3/2}\gamma\bigr)$ --- yakni teorema Lindeberg
        yang sebenarnya dalam bentuk Lyapunovnya.
\end{enumerate}

\medskip
\noindent\textbf{Bagian III --- Dividen kuantitatifnya.}
\begin{enumerate}[resume]
  \item (Fungsi distribusinya) Misalkan $t \in \R$ lalu
        hampirilah $\mathbf 1_{\intoc{-\infty}t}$ dari atas dan
        dari bawah oleh landaian $\mathcal C^3_b$ berlebar $\delta$
        (bangunlah landaiannya, dengan $M_3 = O(\delta^{-3})$).
        Lalu dengan menggabungkannya dengan Bagian I, turunkanlah batas dua sukunya
        \[
        \sup_{t\in\R}\,\bigl|\P(T_n \leq t) -
        \Phi(t)\bigr| \;\leq\;
        \frac{C_1(\beta + \gamma)}{\delta^3\sqrt n} +
        C_2\,\delta
        \qquad (\text{bagi setiap } \delta > 0),
        \]
        dengan konstanta yang eksplisit (sebab suku $C_2\delta$-nya memakai
        bahwa $\Phi$ berkepadatan terbatas oleh
        $\frac1{\sqrt{2\pi}}$), lalu optimumkanlah $\delta \sim
        n^{-1/8}$ untuk memperoleh laju seragam berorde
        $n^{-1/8}$.
        ($n^{-1/2}$ yang optimum --- yakni Berry--Esseen ---
        memerlukan perkakas yang lebih halus; sedangkan intinya adalah
        laju yang \emph{eksplisit} dari penukaran yang elementer.)
  \item (De Moivre--Laplace, yang dikuantifikasi) Khususkanlah pada
        $X_i = 2B_i - 1$ (yakni tanda koin adil): lalu bandingkanlah
        kesimpulannya dengan taksiran lokal
        \cref{pb:b3:product:1}, pertanyaan 7 --- apa yang diberikan
        masing-masing metode yang tak diberikan yang lain?
  \item (Keuniversalan) Jelaskanlah dalam satu paragraf mengapa
        metode penggantiannya menunjukkan lebih dari teorema limit pusatnya: bahwa
        sembarang statistik berbentuk $\E f(\text{jumlah})$ dengan $f$ yang
        mulus tak peka, pada orde $n^{-1/2}$, terhadap
        \emph{seluruh distribusi} sukunya di luar dua momen
        pertamanya --- yakni ``asas invariansi'' yang
        melandasi hasil keuniversalan modern (matriks acak,
        polinomial acak), dan teorema limit pusatnya adalah
        contoh pertamanya.
\end{enumerate}

\medskip
\noindent\textbf{Bagian IV --- Pemulusan, yang didorong: laju
yang lebih baik.} Kehilangan dari $n^{-1/2}$ (bagi $f$ yang mulus) ke
$n^{-1/8}$ (bagi fungsi distribusinya) berasal dari membebankan
$f'''$ dalam norma sup. Padahal hibridanya dapat memperbaiki sebagiannya: sebab
di dalamnya ada suku \omterm{def:b3:clt:gaussianvector}{Gauss}, dan \omterm{def:b3:clt:gaussianvector}{Gauss} \emph{memuluskan}.
\begin{enumerate}[resume]
  \item (\omterm{def:b3:clt:gaussianvector}{Gauss} yang tersembunyi) Untuk $1 \leq i \leq n - 1$, dengan $h =
        \frac{X_i}{\sqrt n}$ atau $\frac{N_i}{\sqrt n}$, dan
        $\theta \in \intcc01$, tulislah $W_i + \theta h = A +
        Z$ dengan $Z = \frac{N_{i+1} + \dots + N_n}{\sqrt
        n}$. Tunjukkanlah bahwa $Z \sim \mathcal N\bigl(0,
        \frac{n-i}n\bigr)$ bebas dari pasangan $(A,
        h)$, lalu simpulkanlah, bagi setiap $g \in
        L^1(\R)$ yang \omterm{def:b3:topology:continuity}{kontinu},
        \[
        \E\bigl[\abs h^3\,\abs{g(W_i + \theta h)}\bigr]
        \;\leq\; \sqrt{\frac{n}{2\pi(n - i)}}\;
        \norm{g}_{L^1}\;\E\abs h^3 .
        \]
  \item Gabungkanlah pertanyaan 10 dengan bentuk integral
        sisa Taylornya,
        \[
        f(w + h) = f(w) + f'(w)h + \tfrac12f''(w)h^2 +
        \int_0^1\frac{(1 - \theta)^2}2\,f'''(w + \theta
        h)\,h^3\,\dd\theta,
        \]
        untuk mengulangi pertanyaan 3--4: bahwa bagi $f \in \mathcal C^3_b$
        yang terlebih lagi $f''' \in L^1(\R)$,
        \[
        \bigl|\E f(T_n) - \E f(G_n)\bigr| \leq
        \frac{\beta + \gamma}{3\sqrt{2\pi}}\cdot
        \frac{\norm{f'''}_{L^1}}{\sqrt n} +
        \frac{M_3(\beta + \gamma)}{6\,n^{3/2}}
        \]
        \emph{(pertanyaan 10 menangani penukaran $i \leq n - 1$
        --- pakailah $\sum_{m=1}^{n-1}m^{-1/2} \leq 2\sqrt n$
        --- sedangkan batas kasar pertanyaan 3 menangani yang
        terakhir)}. Periksalah bahwa landaian pertanyaan 7 memenuhi
        $\norm{\psi_\delta'''}_{L^1} = K_1\delta^{-2}$
        sedangkan $M_3 = K\delta^{-3}$, masukkanlah keduanya, lalu
        optimumkanlah $\delta$: sehingga laju seragam
        fungsi distribusinya membaik menjadi
        $O(n^{-1/6})$.
  \item (Mencocokkan satu momen lagi) Andaikanlah pula $\E
        X_1^3 = 0$ dan $\beta_4 = \E X_1^4 < \infty$.
        Hitunglah $\E N_1^3$ dan $\E N_1^4$, uraikan sampai orde
        keempat, lalu buktikanlah dengan cara yang sama bahwa
        laju fungsi distribusinya menjadi $O(n^{-1/4})$
        \emph{(kini $\norm{\psi_\delta^{(4)}}_{L^1} =
        K_2\delta^{-3}$ dan $M_4 = K'\delta^{-4}$; pilihlah
        $\delta = n^{-1/4}$)}.
  \item (Halangannya) Andaikanlah $k$ momen pertama
        $X_1$ cocok dengan momen Gaussnya ($k = 2$ selalu;
        $k = 3$ tepat saat $\E X_1^3 = 0$; sedangkan $k \geq 4$
        pada dasarnya tak pernah, sebab $\E N_1^4 = 3$). Periksalah bahwa
        skema pertanyaan 10--12 memberikan
        laju fungsi distribusinya
        $n^{-(k-1)/(2k+2)}$, lewat menyeimbangkan
        $\delta^{-k}n^{-(k-1)/2}$ terhadap $\delta$, lalu
        amatilah bahwa eksponennya mendekati
        nilai Berry--Esseen $\frac12$ hanya saat $k \to
        \infty$. Jelaskanlah dalam beberapa kalimat mengapa
        metode penukarannya menjenuh: sebab setiap penukarannya dibebankan dalam
        nilai mutlak, sedangkan jalur Fouriernya
        (yakni ketaksamaan pemulusan Esseen) memanfaatkan
        ayunan selisih \omterm{def:b3:clt:cf}{fungsi karakteristiknya}
        lalu mencapai $C\beta n^{-1/2}$ dengan tiga
        momen saja.
\end{enumerate}

\medskip
\noindent\textbf{Bagian V --- Dua dimensi: teorema limit pusat
berdimensi banyak, lewat penukaran.} Kini misalkan $X_i$ merupakan
\emph{vektor} acak i.i.d.\ yang terpusat di $\R^2$ dengan
matriks kovarians $\Sigma$ dan $\beta' = \E\norm{X_1}^3 <
\infty$ (menurut norma Euclidnya).
\begin{enumerate}[resume]
  \item (\omterm{def:b3:clt:gaussianvector}{Vektor Gauss}, sesuai pesanan) Diagonalkanlah $\Sigma =
        PDP^{\mathsf T}$ (\cref{exo:b3:submanifolds:8}) lalu
        tetapkan $C = P\sqrt DP^{\mathsf T}$. Lalu bagi $Z = (Z^1,
        Z^2)$ sepasang variabel \omterm{def:b3:clt:gaussianvector}{Gauss} baku yang bebas
        (\cref{thm:b3:probability:existence}), tunjukkanlah bahwa $N
        = CZ$ merupakan \omterm{def:b3:clt:gaussianvector}{vektor Gauss}
        (\cref{def:b3:clt:gaussianvector}) berrata-rata $0$,
        berkovarians $\Sigma$, dengan $\gamma' = \E\norm N^3 <
        \infty$; dan bahwa $G_n = \frac{N_1 + \dots +
        N_n}{\sqrt n}$ berdistribusi $\mathcal N(0, \Sigma)$
        \emph{secara persis} bagi salinan $N_i$ yang i.i.d.
  \item (Taylor dalam dua variabel) Untuk $f \colon \R^2 \to
        \R$ berkelas $\mathcal C^3$ dengan $M_3 =
        \max_{\abs\alpha = 3}\sup\abs{\partial^\alpha f} <
        \infty$, buktikanlah
        \[
        \Bigl|f(w + h) - f(w) - \langle\nabla f(w),
        h\rangle - \tfrac12\langle h, D^2f(w)\,h\rangle
        \Bigr| \leq \frac{M_3}6\,\bigl(\abs{h_1} +
        \abs{h_2}\bigr)^3 \leq \frac{\sqrt2\,M_3}3\,
        \norm h^3
        \]
        \emph{(pelajarilah $t \mapsto f(w + th)$ pada
        $\intcc01$)}.
  \item (Teorema limit pusatnya di $\R^2$) Jalankanlah skema penggantiannya pada
        hibrida vektornya $H_i$: tunjukkanlah bahwa suku orde pertama dan
        keduanya meniadakan diri (sebab rata-rata dan kovariansnya cocok), teleskopkanlah,
        lalu naikkanlah seperti pada pertanyaan 5
        (dengan keketatan dari $\E\norm{T_n}^2 =
        \operatorname{tr}\Sigma$; dan pemulusannya kini di
        $\R^2$, \cref{thm:b3:lp:regularization}) untuk
        menyimpulkan: bahwa bagi setiap $f \colon
        \R^2 \to \R$ yang \omterm{def:b3:topology:continuity}{kontinu} terbatas,
        \[
        \E\,f\Bigl(\frac{X_1 + \dots + X_n}{\sqrt n}\Bigr)
        \longrightarrow \E\,f(N), \qquad N \sim \mathcal
        N(0, \Sigma) :
        \]
        yakni \cref{thm:b3:clt:multiclt} pada dimensi $2$, dengan
        laju bagi $f$ yang mulus dan tanpa analisis Fourier.
  \item (Cramér--Wold, dan sebuah fluktuasi bersama) Simpulkanlah
        bahwa $\langle t, \frac{S_n}{\sqrt n}\rangle
        \Rightarrow \mathcal N(0, t^{\mathsf T}\Sigma t)$
        bagi setiap $t \in \R^2$ yang tetap. Penerapannya: bagi
        $(\xi_i)$ real i.i.d.\ yang terpusat, $\E\xi_1^2 = 1$,
        $\E\xi_1^6 < \infty$ (sehingga Bagian V berlaku bagi
        $V_i = (\xi_i, \xi_i^2 - 1)$), tunjukkanlah
        \[
        \frac1{\sqrt n}\Bigl(\sum_{i\leq n}\xi_i,\
        \sum_{i\leq n}(\xi_i^2 - 1)\Bigr) \Longrightarrow
        \mathcal N\Bigl(0, \begin{pmatrix} 1 & \E\xi_1^3\\
        \E\xi_1^3 & \E\xi_1^4 - 1\end{pmatrix}\Bigr) :
        \]
        yakni rata-rata empiris dan momen kedua empirisnya
        berfluktuasi bersama secara \omterm{def:b3:clt:gaussianvector}{Gauss} --- dan bebas pada
        limitnya jika dan hanya jika $\E\xi_1^3 = 0$
        (\cref{thm:b3:clt:gaussianvector}).
\end{enumerate}

\medskip
\noindent\textbf{Bagian VI --- Metode delta.}
\begin{enumerate}[resume]
  \item Misalkan $(\hat\theta_n)$ \omterm{def:b3:probability:space}{peubah acak} dengan
        $\sqrt n(\hat\theta_n - \theta) \Rightarrow
        \mathcal N(0, \sigma^2)$ bagi parameter real
        $\theta$, dan misalkan $g$ terdiferensialkan di
        $\theta$. Buktikanlah \emph{metode delta}-nya:
        \[
        \sqrt n\bigl(g(\hat\theta_n) - g(\theta)\bigr)
        \Longrightarrow \mathcal N\bigl(0,
        g'(\theta)^2\sigma^2\bigr)
        \]
        \emph{(tulislah $g(x) - g(\theta) = (g'(\theta) +
        \eta(x))(x - \theta)$ dengan $\eta \to 0$ di
        $\theta$; tunjukkanlah $\hat\theta_n \to \theta$, lalu
        $\eta(\hat\theta_n) \to 0$, dalam peluang; lalu selesaikanlah
        dengan Slutsky, \cref{exo:b3:clt:8}, dan
        \cref{exo:b3:clt:4}(b))}.
  \item Penerapannya. (a) Bagi $(\xi_i)$ real i.i.d.\ yang
        berrata-rata $\mu$ dan bervarians $\sigma^2$, dan $\bar X_n =
        \frac1n\sum_{i\leq n}\xi_i$: tunjukkanlah $\sqrt n(\bar
        X_n^2 - \mu^2) \Rightarrow \mathcal N(0,
        4\mu^2\sigma^2)$ bila $\mu \neq 0$, dan bahwa bagi
        $\mu = 0$ pernyataan yang benar hidup pada skala lain:
        yakni $n\bar X_n^2 \Rightarrow \sigma^2N^2$ dengan
        $N \sim \mathcal N(0,1)$ (kenalilah fungsi distribusi
        limitnya). (b) (Penstabilan
        varians) Bagi $\hat p_n$ sebagai frekuensi keberhasilan
        cuplikan $\mathcal B(1, p)$, dengan $p \in \intoo01$:
        tunjukkanlah bahwa $g(p) = \arcsin\sqrt p$ memenuhi
        \[
        \sqrt n\,\bigl(g(\hat p_n) - g(p)\bigr)
        \Longrightarrow \mathcal N\Bigl(0, \frac14\Bigr)
        \]
        \emph{apa pun} $p$-nya --- yakni batang galat asimtotik yang
        bebas dari parameter tak diketahuinya; bandingkanlah dengan
        \cref{ex:b3:clt:confidence}.
\end{enumerate}

\medskip
\noindent\textbf{Bagian VII --- Poisson, lewat metode yang sama:
teorema Le Cam.} Penggantiannya mengenal kelas keuniversalan kedua:
yakni jumlah banyak kejadian \emph{langka} yang bebas. Bagi
distribusi pada $\N$, jarak yang tepat adalah \emph{variasi total},
\[
d_{\mathrm{TV}}(\mu, \nu) = \sup_{A\subseteq\N}\,
\abs{\mu(A) - \nu(A)} .
\]
\begin{enumerate}[resume]
  \item Tunjukkanlah bahwa $d_{\mathrm{TV}}(\mu, \nu) =
        \frac12\sum_{k\geq0}\abs{\mu(\{k\}) -
        \nu(\{k\})}$, lalu buktikanlah batas penggandengannya: bahwa bagi
        \emph{sembarang} pasangan $(X, Y)$ \omterm{def:b3:probability:space}{peubah acak} yang
        berdistribusi $\mu$ dan $\nu$ pada ruang yang sama,
        $d_{\mathrm{TV}}(\mu, \nu) \leq \P(X \neq Y)$.
  \item Hitunglah secara persis, bagi $p \in \intoo01$:
        \[
        d_{\mathrm{TV}}\bigl(\mathcal B(1, p), \mathcal
        P(p)\bigr) = p\bigl(1 - \eu^{-p}\bigr) \leq p^2 .
        \]
  \item (Le Cam, lewat penukaran) Misalkan $X_i \sim \mathcal B(1,
        p_i)$ dan $Y_i \sim \mathcal P(p_i)$, dengan $2n$
        variabelnya bebas; lalu $S = X_1 + \dots + X_n$, dan
        ingatlah $Y_1 + \dots + Y_n \sim \mathcal P(\lambda)$
        dengan $\lambda = \sum_ip_i$ (\cref{exo:b3:clt:1}).
        Tukarlah satu koordinat setiap kali pada hibrida
        bulatnya $H_i = Y_1 + \dots + Y_i + X_{i+1} + \dots
        + X_n$: tunjukkanlah, bagi setiap $A \subseteq \N$,
        \[
        \abs{\P(H_{i-1} \in A) - \P(H_i \in A)} \leq
        d_{\mathrm{TV}}\bigl(\mathcal B(1, p_i), \mathcal
        P(p_i)\bigr),
        \]
        lalu simpulkanlah \emph{ketaksamaan Le Cam}:
        \[
        d_{\mathrm{TV}}\bigl(\text{distribusi } S,\ \mathcal
        P(\lambda)\bigr) \leq \sum_{i=1}^np_i^2 .
        \]
  \item Dividennya. (a) Untuk $p_i = \frac\lambda n$: batasnya
        adalah $\frac{\lambda^2}n$ --- yakni hukum kejadian
        langka (\cref{exo:b3:clt:5}) yang dinaikkan ke laju
        eksplisit, seragam atas semua kejadian, dan sahih
        bagi $p_i$ yang tak sama pula. (b) Sebanyak $500$ surat
        dikirimkan, masing-masing tersesat secara bebas dengan
        peluang $\frac1{500}$: batasilah galat model
        Poisson berparameter $1$, lalu taksirlah
        peluang bahwa tak ada surat yang tersesat. (c) Tutuplah
        soalnya: bandingkanlah kedua kelas keuniversalan
        yang ditemui di sini --- \omterm{def:b3:clt:gaussianvector}{Gauss} (banyak sumbangan kecil yang
        terhampar; dua momen yang cocok; Taylor) dan
        Poisson (banyak sumbangan langka; satu rata-rata
        yang cocok; penggandengan variasi-total yang persis) --- beserta
        satu metode penggantian di balik keduanya.

  \item (Galat relatif dan transformasi log) Misalkan $(X_n)$
        i.i.d., positif, berrata-rata $\mu > 0$, bervarians
        $\sigma^2$, dan $\bar X_n$ rata-rata empirisnya. Tunjukkanlah
        lewat metode delta bahwa
        \[
        \sqrt n\,\bigl(\ln\bar X_n - \ln\mu\bigr)
        \Longrightarrow
        \mathcal N\Bigl(0,\ \frac{\sigma^2}{\mu^2}\Bigr) :
        \]
        yakni parameter asimtotik $\ln\bar X_n$ adalah
        \emph{koefisien variasinya} $\sigma/\mu$ ---
        yakni galat relatif yang bebas skala. Lalu simpulkanlah selang
        kepercayaan $95\%$ bagi $\mu$ yang berbentuk
        perkalian $\bar X_n\cdot\eu^{\pm1.96\,\sigma/(\mu\sqrt
        n)}$, lalu jelaskanlah kapan ia lebih disukai daripada
        yang aditif.
  \item (Momen ketiga mengemudikan galatnya) Untuk $X \sim$
        Bernoulli($p$) yang terpusat, hitunglah $\E\bigl[(X -
        p)^3\bigr] = p(1-p)(1-2p)$. Lalu dengan memakai analisis Bagian IV
        (sebab galat penukarannya dikemudikan momen ketiganya), jelaskanlah
        mengapa hampiran normal $\mathcal B(n, p)$
        bersifat tak setangkup bagi $p \neq \frac12$ --- yakni melampaui
        pada satu sisi dan kurang pada sisi lainnya --- dan mengapa
        $p = \frac12$ menikmati laju momen-cocok yang lebih cepat.
        Periksalah tanda kemencengannya secara numerik pada $\mathcal
        B(20, 0.1)$ terhadap $\mathcal N(2, 1.8)$: bandingkanlah
        $\P(S = 0) = 0.9^{20}$ dengan massa \omterm{def:b3:clt:gaussianvector}{Gauss}
        $\intoo{-\infty}{0.5}$.

\end{enumerate}
\end{problem}
